% METHODOLOGIES POUR L'ELECTROMAGNETISME
\chapter{Méthodologies pour l'\'Electromagnétisme}
\chaptermark{Méthodologies}
\section{Maxwell et l'équation d'onde}
\lipsum
\section{L'optique physique}
\lipsum

\begin{figure}[htp]
    \centering
    \tdplotsetmaincoords{70}{140}
\begin{tikzpicture}[tdplot_main_coords,scale=2.5]
    \clip(0,-1.6,-1) rectangle (0,1.6,1.6);

    %% Définition des différents styles
    \tikzstyle{init} = [black]			% base initiale
    \tikzstyle{prec} = [blue]			% 1ere base intermédiaire
    \tikzstyle{nuta} = [red]			% 2eme base initiale
    \tikzstyle{rotp} = [green]			% base finale
    \tikzstyle{base} = [thick,-stealth]	% Tracé des bases
    \tikzstyle{angle} = [thick,-latex]	% Tracé des arcs pour les angles
    \tikzstyle{circle} = [thin,dashed]	% Tracé des cercles
    \tikzstyle{dummy} = [draw=none,black]

    %% Paramètres géométriques
    \def\ephi{30}	% Angle de precession dessiné
    \def\etheta{45}	% Angle de nutation dessiné
    \def\rang{0.6}	% Rayon utilisé pour tracer les angles

    %% Tracé
    % Repère initial
    \coordinate (O) at (0,0,0);
    \node[anchor=south west] at (O) {$O$};
    \draw[base,init] (O) -- (1,0,0) node[anchor=north east]{$\overrightarrow{x}$};
    \draw[base,init] (O) -- (0,1,0) node[anchor=north west]{$\overrightarrow{y}$};
    \draw[base,init] (O) -- (0,0,1) node[anchor=south]{$\overrightarrow{z}$};

    % Précession
    \tdplotsetrotatedcoords{\ephi}{0}{0}
    \draw[tdplot_rotated_coords,angle,prec,dotted] (O) --(1,0,0) node[anchor=north east]{};
    \tdplotdrawarc[tdplot_rotated_coords,circle,prec]{(0,0,0)}{1}{0}{360}{}{}
    \tdplotdrawarc[tdplot_rotated_coords,angle,prec]{(0,0,0)}{\rang}{-\ephi}{0}{anchor=north east,prec}{$\phi$}

    % Nutation
    \tdplotsetrotatedcoords{\ephi}{90-\etheta}{0}
    \draw[tdplot_rotated_coords,base,nuta] (O) --(0,0,1) node[anchor=south east,init]{$P$};
    \draw[tdplot_rotated_coords,dummy] (O) -- (0,0,0.5) node[anchor=south,nuta]{$R$};
    \tdplotsetrotatedthetaplanecoords{0}
    \tdplotdrawarc[tdplot_rotated_coords,circle,nuta]{(0,0,0)}{1}{0}{360}{}{}
    \tdplotdrawarc[tdplot_rotated_coords,angle,nuta]{(0,0,0)}{\rang}{90-\etheta}{0}{anchor=west,nuta}{$\theta$}
\end{tikzpicture}
    \caption{Définition des angles de site ($\theta$) and de gisement ($\phi$).}
    \label{fig:angles-def}
    \hypertarget{def-phi}{}
    \hypertarget{def-theta}{}
\end{figure}
\nomenclature[V]{$\theta$}{Angle de site.}
\nomenclature[V]{$\phi$}{Angle de gisement.}

\nomenclature[S]{OP}{Optique Physique}
\nomenclature[S]{PTD}{Théorie Physique de la Diffraction (Physical Theory of Diffraction)}%

\section[Formulation intégrale en champ électrique et méthode des éléments frontière]{Formulation intégrale en champ électrique et méthode des éléments frontière\sectionmark{EFIE et BEM}}
\lipsum

\nomenclature[S]{BEM}{Méthode des \'Eléments Frontière (Boundary Element Method)}%
\nomenclature[S]{EFIE}{Formulation Intégrale en Champ \'Electrique (Electric Field Integral Equation)}%
\begin{figure}[htp]
    \centering
    \tdplotsetmaincoords{90}{140}
\begin{tikzpicture}[tdplot_main_coords,scale=2.5]
    \tikzstyle{line} = [thick]
    \tikzstyle{arrow} = [->,>=stealth]
    \tikzstyle{vec} = [thick,->,>={latex[scale=2]}]
    \tikzstyle{lrarrow} = [<->,>=to]

    % Sommets des triangles et origine
    \coordinate (O) at (0,0,0);
    \coordinate (a1) at (0,-.5,1.2);
    \coordinate (a2) at (0,.5,1.2);
    \coordinate (a3p) at (1,0,0.2);
    \coordinate (a3m) at (-1,0,0.2);

    % Coordonées pour placer les légendes
    \coordinate (rhonp) at (1.5,0,.7);
    \coordinate (rhonm) at (-1.5,0,.7);
    \coordinate (Tnp) at (1.5,0,1.1);
    \coordinate (Tnm) at (-1.5,0,1.1);

    % Labels de l'origine et du vecteur position
    \node[anchor=north] at (O) {$O$};
    \node[anchor=south] at (O) {$\rr$};

    % Placement des extremites de la longueur ln et du label
    \node (ln1) [above=0.1 of a1] {};
    \node (ln2) [above=0.1 of a2] {};
    \node[anchor=south] (ln) at ($(ln1)!0.5!(ln2)$) {$l_n$};

    % coordonnes des barycentres des deux triangles
    \coordinate (bp) at (barycentric cs:a1=1,a2=1,a3p=1);
    \coordinate (bm) at (barycentric cs:a1=1,a2=1,a3m=1);

    % Coordonnees pour les fleches des legendes de Tn
    % Un peu au dessus du  barycentre parce que c'est plus joli
    \coordinate (btp) at ($(bp)!0.1!(ln)$);
    \coordinate (btm) at ($(bm)!0.1!(ln)$);

    % Coordonnees pour l'extremites des vecteurs rn
    % Un peu en dessous du barycentre parce que c'est deng shui
    \coordinate (brp) at ($(bp)!0.2!(a3p)$);
    \coordinate (brm) at ($(bm)!0.2!(a3m)$);

    % Coordonnees pour les fleches des legendes de rho_n
    % Au milieu des rn
    \coordinate (rnp) at ($(brp)!0.5!(a3p)$);
    \coordinate (rnm) at ($(brm)!0.5!(a3m)$);

    % ecriture des labels
    \node[anchor=south] at (rhonp) {$\boldsymbol{\rho}_n^+ = \rr-\rr_n^+$};
    \node[anchor=south] at (rhonm) {$\boldsymbol{\rho}_n^- = \rr-\rr_n^-$};

    \node[anchor=east] at (Tnp) {$T_n^+$};
    \node[anchor=west] at (Tnm) {$T_n^-$};

    \node[anchor=north] at (a3p) {$\rr_{n^+}$};
    \node[anchor=north] at (a3m) {$\rr_{n^-}$};

    % Dessing des triangles
    % ATTENTION le triangle de derriere doit etre dessine en premier,
    % sinon il apparaît devant
    \draw[fill=white] (a1) -- (a2) -- (a3m) -- cycle;
    \draw[fill=white] (a1) -- (a2) -- (a3p) -- cycle;

    % Double fleche ln
    \draw[lrarrow] (ln1) -- (ln2);

    % Vecteurs positions
    \draw[dashed,vec] (O) -- (brp);
    \draw[dashed,vec] (O) -- (brm);

    % Vecteurs rn
    \draw[vec] (a3p) -- (brp);
    \draw[vec] (brm) -- (a3m);

    % Fleches des legendes
    \draw[arrow] (rhonp) to[bend right] (rnp);
    \draw[arrow] (rhonm) to[bend left] (rnm);

    \draw[arrow] (Tnp) to[bend left] (btp);
    \draw[arrow] (Tnm) to[bend right] (btm);
\end{tikzpicture}
    \caption{Definition des fonctions de base de Rao-Wilton-Glisson.}
    \label{fig:RWG-def}
\end{figure}
\nomenclature[S]{RWG}{Fonctions de base de Rao-Wilton-Glisson}

\section{Différentiation de la BEM}
\lipsum